Sharpness-aware minimisation (SAM) \cite{foret2020sharpness} trains on the worst-case loss within a small ball around the weights, and has been reported to improve generalisation and, in particular, to be robust to noisy labels. This is usually motivated by the idea that flat minima generalise better \cite{keskar2016large}, yet the evidence that sharpness explains generalisation is mixed \cite{jiang2019fantastic,andriushchenko2023modern}, and recent analyses suggest that SAM's behaviour under label noise is not simply a flatness effect \cite{baek2024why,andriushchenko2022understanding}. Networks that can fit random labels \cite{zhang2016understanding} make label noise a natural stress test: an optimiser must avoid memorising the flipped labels.

We ask three questions at a scale that fits a laptop CPU, with a validation-selected $\rho$ (6 values) for SAM, a validation-selected $\lambda$ (4 values) for SGD+WD, and a shared, untuned learning rate:
\textbf{H1}: does SAM improve clean test accuracy over SGD and over SGD with tuned weight decay under label noise? \textbf{H2}: how sensitive is SAM to the neighbourhood size $\rho$, and is there an interior optimum? \textbf{H3}: does a measured sharpness (loss increase under a fixed-norm weight perturbation) track the generalisation gap?

Our contributions are modest and empirical: (i) a three-optimiser comparison on digits with 0/20/40\% symmetric label noise and on a two-spirals task, with 5 evaluation seeds and hyper-parameter selection on separate tuning seeds; (ii) a $\rho$ sweep, a weight-decay sweep and two ablations (random-direction perturbation, SAM+WD); (iii) a sharpness--gap correlation analysis over 260 runs. Our headline finding is that SAM beats plain SGD by a wide margin under noise but only narrowly, and not robustly, beats SGD with weight decay, and that fixed-radius sharpness does not track the gap across optimisers.
